\documentclass[11pt]{article}
\usepackage[margin=0.82in]{geometry}
\usepackage{amsmath,amssymb,amsthm}
\usepackage{microtype}
\usepackage{xcolor}
\usepackage{hyperref}
\hypersetup{colorlinks=true,linkcolor=blue!55!black,urlcolor=blue!55!black,citecolor=blue!55!black}
\newtheorem{theorem}{Theorem}
\newtheorem{lemma}[theorem]{Lemma}
\newcommand{\sigmatop}{\sigma(X,Y)}
\title{A bounded $\sigma$-closed operator whose entire resolvent\\fails $\sigma$-continuity}
\author{Candidate counterexample packet}
\date{11 August 2026}

\begin{document}
\maketitle

\begin{center}
\fbox{\parbox{0.92\linewidth}{\textbf{Status:} candidate counterexample likely valid; full negative answer to the resolvent question in arXiv:0903.1001.  On an explicit norming dual pair, an everywhere-defined bounded $\sigma$-closed involution $A$ satisfies $\rho(A)=\mathbb C\setminus\{-1,1\}$ but $\rho_\sigma(A)=\varnothing$.}}
\end{center}

\section{Source question}

For a norming dual pair $(X,Y)$, write $\sigma=\sigma(X,Y)$.  If $A$ is
$\sigma$-closed, Kunze defines
\[
 \rho_\sigma(A)=\{\lambda\in\rho(A):R(\lambda,A)
 \text{ is $\sigma$-continuous}\}
\]
and asks whether $\rho_\sigma(A)=\rho(A)$ must always hold
\cite[Section~3]{Kunze}.  The answer is no, even when $A$ is bounded,
everywhere defined, and invertible.

\section{Proof Intuition}

A bounded operator is $\sigma(X,Y)$-continuous exactly when its adjoint
leaves $Y$ invariant.  Closedness of its graph is weaker: it is enough that a
separating core inside $Y$ be invariant under the adjoint.  We exploit this
gap on $X=\ell^1$.

Take $Y$ to contain the usual norming subspace $c_0$ and one extra bounded
sequence $h$.  An adjacent-coordinate permutation preserves $c_0$, so its
graph is detected as closed by the coordinate functionals.  But it sends
$h$ outside $Y$, destroying weak continuity.  Because the permutation is an
involution, every resolvent is an affine combination of the identity and the
same permutation, and the defect persists throughout the resolvent set.

\section{The norming dual pair}

Let
\[
 X=\ell^1(\mathbb N),\qquad
 h=(1,0,1,0,\ldots)\in\ell^\infty,
 \qquad Y=c_0+\mathbb C h\subset\ell^\infty,
\]
with the canonical pairing
$\langle x,y\rangle=\sum_{n\geq1}x_ny_n$ and the inherited norms.

\begin{lemma}\label{lem:dualpair}
The pair $(X,Y)$ is a norming dual pair.  Moreover, $Y$ is a closed proper
subspace of $X^*=\ell^\infty$.
\end{lemma}

\begin{proof}
The space $c_0$ is closed in $\ell^\infty$, and $h\notin c_0$; hence its
one-dimensional extension $Y=c_0+\mathbb Ch$ is closed.  Since $c_0\subset
Y$, finitely supported sign sequences show that
\[
 \|x\|_1=\sup\{ |\langle x,y\rangle|:y\in Y,\ \|y\|_\infty\leq1\}.
\]
Conversely, every $y\in Y\subset\ell^\infty=(\ell^1)^*$ has dual norm
exactly $\|y\|_\infty$.  Thus the two spaces norm each other.  Finally, the
constant sequence $\mathbf1$ does not belong to $Y$: if
$\mathbf1=z+ch$ with $z\in c_0$, the odd coordinates force $c=1$, while the
even coordinates would force $z_{2n}=1$, a contradiction.
\end{proof}

\section{Counterexample}

Define the adjacent-swap isometry $P:X\to X$ by
\[
 (Px)_{2n-1}=x_{2n},\qquad (Px)_{2n}=x_{2n-1}\qquad(n\geq1).
\]
Its Banach adjoint $P^*$ performs the same permutation on $\ell^\infty$.

\begin{theorem}\label{thm:counterexample}
The bounded operator $A=P$ is $\sigma(X,Y)$-closed, but
\[
 \rho(A)=\mathbb C\setminus\{-1,1\},
 \qquad \rho_\sigma(A)=\varnothing.
\]
In particular, $\rho_\sigma(A)\neq\rho(A)$.
\end{theorem}

\begin{proof}
First observe that $P^*c_0=c_0$.  Suppose a net satisfies
$x_\alpha\to x$ and $Px_\alpha\to z$ in $\sigma(X,Y)$.  For every
$y\in c_0$,
\[
 \langle z,y\rangle
 =\lim_\alpha\langle Px_\alpha,y\rangle
 =\lim_\alpha\langle x_\alpha,P^*y\rangle
 =\langle x,P^*y\rangle
 =\langle Px,y\rangle.
\]
Because $c_0$ separates the points of $\ell^1$, $z=Px$.  Thus the graph of
$P$ is $\sigma\times\sigma$ closed.

On the other hand,
\[
 P^*h=(0,1,0,1,\ldots)=\mathbf1-h\notin Y.
\]
Indeed, membership of $\mathbf1-h$ in $c_0+\mathbb Ch$ would imply
membership of $\mathbf1$.  Therefore $P$ is not $\sigma$-continuous, by the
standard adjoint-invariance criterion for norming dual pairs recalled in the
source.

Since $P^2=I$, its spectrum is $\{-1,1\}$ and for
$\lambda\notin\{-1,1\}$,
\[
 R(\lambda,P)=(\lambda I-P)^{-1}
 =\frac{\lambda I+P}{\lambda^2-1}.
\]
Applying the adjoint to $h$ gives
\[
 R(\lambda,P)^*h
 =\frac{\lambda h+P^*h}{\lambda^2-1}
 =\frac{\mathbf1+(\lambda-1)h}{\lambda^2-1}.
\]
This sequence is not in $Y$: all of its even coordinates have the same
nonzero constant value $(\lambda^2-1)^{-1}$, whereas the even coordinates of
every element of $c_0+\mathbb Ch$ converge to zero.  Hence no ordinary
resolvent of $P$ is $\sigma$-continuous, proving the claim.
\end{proof}

\section{Scope and novelty}

The counterexample uses a proper closed norming subspace $Y\subset X^*$.
That is essential to this mechanism: for the full dual pair $(X,X^*)$, every
bounded operator is weakly continuous.  No semigroup, positivity, or
unbounded-domain pathology is involved; failure occurs already for a
surjective linear isometry of order two.

A bounded search through 11 August 2026 covered the run's result, attempt,
and proof-gap indexes and searches for the exact question, the notation
$\rho_\sigma(A)$, $\sigma$-closed resolvents, and norming dual pairs.  It found
the original question and subsequent semigroup literature, but no later
resolution or matching counterexample.  Novelty confidence is moderate
pending specialist review; the construction itself is elementary and fully
explicit.

\begin{thebibliography}{9}
\bibitem{Kunze}
M. Kunze,
\emph{Continuity and equicontinuity of transition semigroups},
Semigroup Forum 79 (2009), 540--560; arXiv:0903.1001.
\url{https://arxiv.org/abs/0903.1001}.
\end{thebibliography}

\end{document}
